% TOP_PROBLEM: {"id": 1463, "title": "Stable Field Conjecture", "category_id": 4, "category": {"id": 4, "name": "algebra", "display_name": "Algebra", "description": "Group theory, ring theory, field theory, and algebraic structures.", "slug": "algebra", "order_index": 4, "created_at": "2026-07-31T15:26:25.671Z"}, "status": "open"}
% FIRST_POSTED: 2026-10-01
% ATTEMPT_STATUS: unresolved
% Imported from: unsolvedmath_additional_500.tex; UnsolvedMath ID 1463
% !TeX program = lualatex
% Compile twice: lualatex 1463.tex
% Self-contained: no external bibliography, images, or \input files.
% Numeric section labels and visible numbers are original UnsolvedMath IDs.
\documentclass[11pt,a4paper]{article}
\usepackage[margin=24mm,headheight=14pt]{geometry}
\usepackage{fontspec}
\setmainfont{Latin Modern Roman}
\setsansfont{Latin Modern Sans}
\setmonofont{Latin Modern Mono}
\usepackage{amsmath,amssymb,mathtools}
\usepackage{unicode-math}
\setmathfont{Latin Modern Math}
\usepackage{microtype}
\usepackage{enumitem,xurl,needspace}
\usepackage[dvipsnames]{xcolor}
\usepackage{hyperref}
\hypersetup{unicode=true,colorlinks=true,linkcolor=MidnightBlue,urlcolor=MidnightBlue,
 pdftitle={UnsolvedMath Problem 1463},pdfauthor={Source-annotated compilation},bookmarksnumbered=true}
\usepackage{bookmark,tocloft,titlesec,fancyhdr}
\setlength{\cftsecnumwidth}{4.2em}
\cftsetpnumwidth{2.5em}
\cftsetrmarg{4em}
\titleformat{\section}{\Large\bfseries\raggedright}{\thesection}{0.7em}{}
\pagestyle{fancy}\fancyhf{}
\fancyhead[L]{\small\sffamily Additional UnsolvedMath statements}
\fancyhead[R]{\small Original catalogue IDs}
\fancyfoot[C]{\thepage}
\setlength{\parindent}{0pt}
\setlength{\parskip}{5pt plus 1pt}
\setlength{\emergencystretch}{3em}
\setcounter{tocdepth}{1}\setcounter{secnumdepth}{1}
\setlist[itemize]{leftmargin=3em,itemsep=4pt,topsep=4pt,parsep=0pt}
\allowdisplaybreaks
\newcommand{\N}{\mathbb{N}}
\newcommand{\Z}{\mathbb{Z}}
\newcommand{\Q}{\mathbb{Q}}
\newcommand{\R}{\mathbb{R}}
\newcommand{\C}{\mathbb{C}}
\newcommand{\F}{\mathbb{F}}
\newcommand{\A}{\mathbb{A}}
\newcommand{\PP}{\mathbb{P}}
\newcommand{\E}{\mathbb{E}}
\newcommand{\eps}{\varepsilon}
\newcommand{\dd}{\,\mathrm d}
\newcommand{\sref}[2]{\hyperlink{source:#1:#2}{[S#2]}}
\newif\ifshowcatalogue
\showcataloguetrue
\newcommand{\cataloguescope}[1]{\ifshowcatalogue
 \paragraph{Statement recorded in the supplied source.}
 #1\fi}
\begin{document}
% UnsolvedMath ID 1463; source catalogue code MOD-006

\setcounter{section}{1462}

\section{The Stable Field Conjecture}\label{1463}

{\small\sffamily UnsolvedMath ID 1463 · Algebra}

\subsection{Definitions and mathematical statement}

\paragraph{Shared notation.}Unless a section says otherwise, $\N=\{1,2,\ldots\}$,
$\Z$ is the ring of integers, $\Q,\R,\C$ are the rational, real, and complex
fields, $[N]=\{1,\ldots,\lfloor N\rfloor\}$, and $\log$ is the natural
logarithm. For real functions of a parameter tending to infinity,
$f\ll g$ and $f=O(g)$ mean $|f|\leq Cg$ eventually for some fixed $C$;
$f\gg g$ reverses this comparison; $f=o(g)$ means $f/g\to0$;
$f\sim g$ means $f/g\to1$; and $f\asymp g$ means both order bounds.
A subscript on $O$ or $\ll$ specifies allowed constant dependence.
The expression $n^{o(1)}$ in an upper bound means growth smaller than
$n^\epsilon$ for each fixed $\epsilon>0$. Local definitions govern any
exception, finite-field convention, or use of the same symbol for another
object. An asymptotic claim is not automatically an effective algorithm.



Use the pure first-order language of fields. A theory is stable when no formula $\varphi(x;y)$ has the order property: there are no sequences $(a_i),(b_j)$ in a model with $\varphi(a_i;b_j)$ holding exactly when $i<j$. A field is separably closed if it has no proper finite separable algebraic extension. The assertion concerns all infinite fields with stable theory. \sref{1463}{1} \sref{1463}{2}

\paragraph{Statement recorded in the supplied source.}
Is every infinite field with a stable first-order theory separably closed? \sref{1463}{1}

\paragraph{Residual task reported by the archive.}

The embedded review (2026-08-17) records: Show every infinite stable field is separably closed. \sref{1463}{1}

\subsection{Short English statement}

Must an infinite field whose first-order theory is stable already be separably closed?

\subsection{Research attempt: explicit order-property witnesses from orders and valuations}
\textbf{Status.} We rule out two important mechanisms for a non-separably-closed field to be stable, but do not show those mechanisms cover every field. The primary research paper [R1] concerns the stable-fields conjecture and the étale-open topology. We use it for context, not as a theorem that every stable field is already known to be separably closed.

\subsubsection{1. A definable infinite order destroys stability}
Suppose a definable relation orders an infinite sequence $a_0<a_1<\cdots$ in a field. The same defining formula, with any fixed defining parameters included as additional parameters, has the order property on $(a_i),(a_j)$. This contradicts stability. It is enough to build one such formula; no decision procedure for the entire field theory is required.

For a real closed field the order is definable in the pure field language:
\[x<y\quad\Longleftrightarrow\quad
\exists z\;(z\ne0\ \wedge\ z^2=y-x).
\]
Indeed the nonzero squares are exactly the positive elements. Taking $a_i=i\cdot1$ and $b_j=j\cdot1$ gives the order pattern exactly when $i<j$. Thus real closed fields are unstable even when no order symbol is included in the language. Merely removing a relation symbol cannot remove an order that is already definable.

\subsubsection{2. A definable nontrivial valuation also supplies the pattern}
Suppose a nontrivial valuation ring $\mathcal O\subset K$ is definable with parameters. Choose $t\in K^*$ with positive valuation. Such an element exists by nontriviality, replacing an element by its inverse if necessary. For nonzero $x,y$, define
\[\psi(x;y)\quad\Longleftrightarrow\quad
\neg\exists z\;(x=yz\ \wedge\ z\in\mathcal O).
\]
In the field there is exactly one solution $z=x/y$, so
\[\psi(t^i;t^j)\quad\Longleftrightarrow\quad
v(t^i/t^j)=(i-j)v(t)<0
\quad\Longleftrightarrow\quad i<j.
\]
The powers are distinct because their valuations are distinct. This gives the order property without assuming the value group is Archimedean or that the valuation is discrete. Parameters defining $\mathcal O$ do not rescue stability, since they can be fixed throughout the array. Therefore a stable field has no definable nontrivial valuation ring.

\subsubsection{3. Why the algebraic target is still stronger}
Failure of separable closure supplies an irreducible separable polynomial of degree greater than one. Such a polynomial does not on its own define either a linear order or a nontrivial valuation ring. Turning that algebraic failure into one of the preceding formulas would be an additional theorem. Conversely, the absence of these two definable objects is not a characterization of stability: other formulas can have the order property.

One must also distinguish a valuation existing on the abstract field from a valuation ring definable in its pure field structure. The formula above uses membership in $\mathcal O$ explicitly. Naming an arbitrary external valuation changes the structure and can change its stability, so instability of the expansion alone does not refute stability of the original field.

\subsubsection{4. Verification and residual problem}
The algebraic checks are exact: the square formula uses real closedness, the valuation formula uses division only for nonzero parameters, and the sign computation uses positivity of $v(t)$ in an ordered abelian group. The diagnostic checks the finite order arrays for indices 0 through 20 symbolically; the formula proves arbitrarily large arrays and thus the full order property. No numerical experiment asserts stability or separable closure. The missing step is a theorem that an arbitrary infinite stable field with a proper finite separable extension must generate some order-property configuration. The explicit witnesses above settle only the stated ordered and definably valued subcases.

\subsection{Sources}

{\small

\begin{itemize}

\item[\hypertarget{source:1463:1}{S1}] ULAM.AI, \emph{UnsolvedMath}, supplied \texttt{problems.json}, record 1463 (MOD-006), statement and background. The embedded literature-review date is 2026-08-17; that is the archive’s review date, not a fresh certification. Dataset landing page: \url{https://huggingface.co/datasets/ulamai/UnsolvedMath}.

\item[\hypertarget{source:1463:2}{S2}] Finite Undecidability in NIP Fields, J. Symbolic Logic (2023). \url{https://doi.org/10.1017/jsl.2022.59}. Bibliographic information imported from the supplied record.

\item[R1] The étale-open topology and the stable fields conjecture, Journal of the European Mathematical Society 26 (2024), 4033--4070, primary research paper.
\url{https://ems.press/content/serial-article-files/48118?nt=1}.
\end{itemize}

}

\label{end:1463}
\end{document}
